% Exercise sheet for the journal club talk on Gardiner, Chapter 5 (../fpe-slides/FPE_ch5.tex).
% The solutions are in wright-fisher-key.tex. Equation numbers are Gardiner, 2nd ed.
\documentclass[11pt]{article}

\usepackage[letterpaper,margin=1in]{geometry}
\usepackage[T1]{fontenc}
\usepackage{lmodern}
\usepackage{microtype}
\usepackage{amsmath,amssymb,mathtools}
\usepackage[dvipsnames]{xcolor}
\usepackage{enumitem}
\usepackage[colorlinks=true,urlcolor=PineGreen]{hyperref}

\setlength{\parindent}{0pt}
\setlength{\parskip}{6pt}
\setlist{itemsep=3pt,topsep=2pt}

\newcommand{\dd}{\mathrm{d}}
\newcommand{\avg}[1]{\langle #1\rangle}

% Title and running header at the top of each page
\newcommand{\sheethead}[1]{%
  {\sffamily\Large\color{PineGreen}Genetic drift, with and without mutation}\par\vspace{2pt}
  {\sffamily\small Exercise sheet\hfill #1}\par\vspace{-7pt}
  \rule{\textwidth}{0.5pt}\par\vspace{2pt}}
\newcommand{\exsec}[1]{\par\vspace{4pt}{\sffamily\bfseries\large\color{PineGreen}#1}\par\vspace{-2pt}}

\begin{document}

% =========================================================
\sheethead{Journal club}

\textbf{Model.} A population carries $N$ copies of a gene, each of type $A$ or type $a$, and $x$ is the fraction
of type $A$. Neither type is fitter. Each generation,
\begin{enumerate}
  \item every copy mutates to the other type, independently, with probability $u$;
  \item the next generation is $N$ copies drawn at random, with replacement, from the mutated pool
        (Wright--Fisher resampling).
\end{enumerate}
Time $t$ is counted in generations. For large $N$, with $Nu$ held fixed, $x(t)$ is close to a diffusion on
$0<x<1$. Throughout,
\[ \theta=2Nu. \]
Beware a clash of names: biologists call the randomness of resampling \emph{genetic drift}. It is the noise
$B$, not the drift $A$.

\textbf{Convention:} as in the talk, $B$ is the variance per unit time, $\avg{(\Delta x)^2}=B\,\Delta t$. Parts
(a)--(d) are the core exercise; part (e) is an optional extension. Equation numbers refer to Gardiner, 2nd edition.

\exsec{(a) Build the forward equation}
Given $x$, find the mean and the mean square of the change $\Delta x$ over one generation. Keep the leading order
in $1/N$, counting $u$ as order $1/N$. Show that
\[ A(x)=u\,(1-2x),\qquad B(x)=\frac{x(1-x)}N, \]
and write the forward equation for $p(x,t)$ and the backward operator $\mathcal L^\dagger$.
\begin{itemize}
  \item Which term comes from mutation, and which from resampling?
  \item Why does $B$ vanish at $x=0$ and $x=1$, and why does $A$ not?
  \item Roughly how many generations does it take for $x$ to change by an amount of order one?
\end{itemize}
\emph{Hint:} After mutation the fraction is $x'=x(1-u)+(1-x)u$. Resampling makes $N$ times the new fraction a
binomial variable, with $N$ trials and success probability $x'$.

\vfill
\newpage

% =========================================================
\sheethead{Core exercise, continued}

\exsec{(b) Find the stationary state}
\begin{enumerate}
  \item Use the potential solution (5.2.39) to show that
        \[ p_s(x)=\frac{\Gamma(2\theta)}{\Gamma(\theta)^2}\,\bigl[x(1-x)\bigr]^{\theta-1}. \]
        Sketch it for $\theta=\tfrac14$, $1$ and $4$. What competes with what? Why is $p_s$ U-shaped when
        mutation is rare?
  \item Classify the boundary $x=0$ with Feller's test (5.2.25)--(5.2.28). Show that it is regular for
        $0<\theta<1$ and an entrance boundary for $\theta\ge1$.
  \item When $x=0$ is regular, a boundary condition has to be chosen. Which one does the discrete model choose,
        and why? Why does that make the potential solution the right one?
  \item What does Gardiner's quick rule (5.2.20)--(5.2.24) predict at $x=0$? Why does it fail? Compare the
        squared-radius example from the talk, $\dd y=(dD-2ky)\,\dd t+\sqrt{4Dy}\,\dd W$: near $x=0$, what
        plays the role of the dimension $d$?
\end{enumerate}
\emph{Hints:} Near $x=0$, keep only the leading power of $x$ in $A$, in $B$, and in Feller's functions $f$ and $g$.
The normalisation is a Beta integral, $\int_0^1[x(1-x)]^{\theta-1}\,\dd x=\Gamma(\theta)^2/\Gamma(2\theta)$
(\S\S5.2.1, 5.2.2).

\vfill
\newpage

% =========================================================
\sheethead{Core exercise, continued}

\exsec{(c) Find the eigenfunctions}
The backward operator is $\mathcal L^\dagger f=u(1-2x)\,f'+\dfrac{x(1-x)}{2N}\,f''$.
\begin{enumerate}
  \item Show that $\mathcal L^\dagger$ maps a polynomial of degree $n$ to a polynomial of degree $n$. Deduce that
        each eigenfunction $Q_n$ is a polynomial of degree $n$, and read off from the coefficient of $x^n$ that
        \[ \lambda_n=\frac{n(n-1+2\theta)}{2N},\qquad n=0,1,2,\dots \]
  \item Find $Q_1$ and $Q_2$.
  \item Show that $\avg{x(t)}-\tfrac12=(x_0-\tfrac12)\,e^{-2ut}$ exactly. Why does $N$ not appear? Which rate is
        the spectral gap?
  \item Here $P_n=p_s\,Q_n$. Why? The heterozygosity $H=2x(1-x)$ is the probability that two copies drawn at random
        differ. Find its stationary average,
        \[ H_s=\frac\theta{1+2\theta}, \]
        without doing an integral.
  \item Interpret the two terms of $\lambda_2=1/N+4u$.
\end{enumerate}
\emph{Hints:} Try $Q_2=x(1-x)+c$. The relation $P_\lambda=p_sQ_\lambda$ comes from (5.2.71)--(5.2.76). Since $Q_0=1$,
biorthogonality says $\int P_n\,\dd x=0$ for $n\ge1$ (5.2.83).

\vfill
\newpage

% =========================================================
\sheethead{Core exercise, continued}

\exsec{(d) Switch off mutation}
Now set $u=0$, so $A=0$ and $B=x(1-x)/N$. Label the modes by their degree, as in (c), so here the absorption rate
is $\lambda_2$ and the rate of approach to the QSD is $\lambda_3-\lambda_2$ (in the talk, $\lambda_1$ and
$\lambda_2-\lambda_1$).
\begin{enumerate}
  \item Show that there is no stationary density. What type are the ends now? What happens to the population?
  \item Now $\lambda_n=n(n-1)/2N$, so $\lambda_0=\lambda_1=0$. Using $Q_1=x$, show that $\avg{x(t)}=x_0$ for all
        $t$. Deduce the probability that type $A$ takes over the whole population, and check it against the
        splitting probability (5.2.186)--(5.2.190).
  \item $Q_2=x(1-x)$ decays at $\lambda_2=1/N$. Show that $P=Q/[x(1-x)]$ solves the forward equation whenever $Q$
        solves the backward one. Compare $P_\lambda=p_sQ_\lambda$: what plays the role of $p_s$? Deduce that
        \emph{the QSD is uniform on $(0,1)$}. Normalise so that $\int P_2Q_2\,\dd x=1$, and show that the
        probability of keeping both types is $S(t)\simeq6x_0(1-x_0)\,e^{-t/N}$ at long times.
  \item Unlike the $\sin(\pi x/L)$ QSD in the talk, this one does not vanish at the absorbing ends. Find the current
        out through each end, and check that the total is $\lambda_2$. Why can $p$ stay finite at an end here?
  \item Check that $Q_3=x(1-x)(1-2x)$, and find $\lambda_3$. At what rate does $p/S$ approach the QSD? What changes
        if $x_0=\frac12$?
  \item Find the mean time $T(x_0)$ until one type is lost, from (5.2.154), and evaluate $T(\frac12)$. Explain why
        the average of $T$ over the QSD must be $1/\lambda_2=N$. Check it.
\end{enumerate}
\emph{Hints:} With $A=0$, (5.2.154) reads $\tfrac12BT''=-1$. You may need $\int_0^1x\ln x\,\dd x=-\tfrac14$.
What happens to $p(x,t)$ if it starts as the QSD? (\S\S5.2.5, 5.2.7, 5.2.8)

\vfill
\newpage

% =========================================================
\sheethead{An optional extension}

\exsec{(e) Optional: from one to the other}
Let mutation be rare, $\theta\to0$ at fixed $N$.
\begin{enumerate}
  \item Show that $p_s\to\tfrac12\delta(x)+\tfrac12\delta(x-1)$, in the sense that a fixed neighbourhood of each end
        holds probability $\tfrac12$. Show that $\lambda_1\to0$. Which conserved quantity of (d) does $Q_1=1-2x$
        become?
  \item Show that
        \[ \frac{p_s\,Q_2}{\theta/2}\;\longrightarrow\;1-\tfrac12\delta(x)-\tfrac12\delta(x-1). \]
        Its interior is the QSD of (d). What do the $\delta$ terms mean? Check that this limit is orthogonal to
        both zero modes of (d), $Q_0=1$ and $Q_1=x$.
  \item When $\theta\ll1$, the population almost always consists of a single type. New mutants appear at rate $Nu$
        per generation, and by (d) each takes over with probability $1/N$. Show that this two-state picture
        relaxes at exactly $\lambda_1$. So neutral variants replace each other at the mutation rate, whatever the
        population size: Kimura's result, the basis of the molecular clock. Compare the double well in the talk.
        Which two time scales must be well separated, and when does the picture fail?
\end{enumerate}
\emph{Hints:} As $\theta\to0$, $\Gamma(2\theta)/\Gamma(\theta)^2\simeq\theta/2$, and
$\int_0^\varepsilon x^{\theta-1}\,\dd x=\varepsilon^\theta/\theta$. Hold $\varepsilon$ fixed while $\theta\to0$.

\vfill

{\footnotesize The Wright--Fisher diffusion is additional material beyond the book examples. References:
M.~Kimura, \emph{Solution of a process of random genetic drift with a continuous model},
\href{https://doi.org/10.1073/pnas.41.3.144}{Proc.\ Natl.\ Acad.\ Sci.\ USA \textbf{41}, 144 (1955)}, which solves
part (d); M.~Kimura, \emph{Diffusion models in population genetics},
\href{https://doi.org/10.1017/S0021900200108368}{J.\ Appl.\ Probab.\ \textbf{1}, 177 (1964)}; M.~Kimura,
\emph{Evolutionary rate at the molecular level},
\href{https://doi.org/10.1038/217624a0}{Nature \textbf{217}, 624 (1968)}, for part (e). The accompanying key uses
the same part labels.\par}

\end{document}
